\documentclass[12pt,a4paper]{article}

\usepackage{german}
\usepackage[utf8]{inputenc}
\usepackage{amsfonts}
\usepackage{amsthm}
\usepackage{amssymb}
\usepackage{enumerate}

\usepackage{graphicx}
\usepackage{float}
%\usepackage{floatflt}
\usepackage{subfigure}

\oddsidemargin=0cm \headheight=-2.1cm
\textwidth=16.4cm \textheight=26cm

\newcommand{\wi}{\angle}
\newcommand{\bogen}{\widehat}
\newcommand{\grad}{^\circ}
\newcommand{\de}{\triangle}
\newcommand{\ab}{\hspace{5pt}}\parindent0pt

\renewcommand{\figurename}{Abb.}

\newcounter{fort}[subsection]
\newcounter{fort2}[subsection]
\newtheorem{bsp}[fort]{Beispiel}
\newtheorem{satz}[fort2]{Satz}

\begin{document}
\thispagestyle{empty}
\begin{figure}[h]
\includegraphics[height=3.1cm]{logo}
\end{figure}

\vspace{1cm}

\begin{center}
\Huge{\textbf{Vorrundenprüfung SMO 2014}}\\[1.5cm]
\large{Lugano, Lausanne, Zürich - 11. Januar 2014}\\[3.5cm]
\end{center}


Bitte beachtet die folgenden Punkte:

\begin{itemize}
\item Da eure Lösungen kopiert und nach Aufgaben sortiert werden, sollen alle Blätter nur auf einer Seite beschrieben werden und für jede Aufgabe ist ein neues Blatt anzufangen. Um alles, was uns das Korrigieren erleichtert, sind wir sehr froh, hebt z.B. wichtige Stellen hervor oder markiert Überflüssiges.

\item Schreibt ALLES auf, was ihr bei einer Aufgabe herausfindet, auch scheinbar triviale Beobachtungen können Teilpunkte geben. Und gib ALLE lesbaren Blätter ab, selbst für falsche oder widersprüchliche Aussagen werden niemals Punkte abgezogen.

\item Die Aufgaben sind nach Schwierigkeit geordnet, es kommt aber sicherlich vor, dass jemand den Kniff einer einfachen Aufgabe nicht sieht, dafür eine der schwierigen Probleme perfekt löst. Lasst euch also trotzdem nicht von der Nummer einer Aufgabe beeindrucken. Es geben alle Aufgaben gleich viele Punkte. Ihr habt genügend Zeit, euch mit jeder Aufgabe eine Zeit lang auseinanderzusetzen, wir raten euch unbedingt dies zu tun.

\item Man kann die Prüfung vorzeitig abgeben und den Raum leise verlassen.

\item Damit uns beim Korrigieren nichts verloren geht, soll am Schluss auf jedem Blatt euer Name (oder mindestens ein Kürzel), die Aufgabennummer, die Seitenzahl und die totale Anzahl Seiten zu dieser Aufgabe stehen. Also z.B. "`C.F.Gauss Nr. 3 S. 2/4"'.

\item Deine korrigierte Prüfung wird dir in der nachfolgenden Woche mit beiliegendem Couvert zugeschickt. Bitte adressiere es jetzt korrekt, gross und gut leserlich mit deiner Adresse.
\end{itemize}

\end{document}
